\documentclass[10pt,letterpaper]{article}
\usepackage[T1]{fontenc}
\usepackage[utf8]{inputenc}
\usepackage{lmodern}
\usepackage[margin=0.76in]{geometry}
\usepackage[table]{xcolor}
\usepackage{array,tabularx,longtable,booktabs}
\usepackage{enumitem}
\usepackage{amssymb,amsmath}
\usepackage{hyperref}
\usepackage{fancyhdr,lastpage}
\usepackage{microtype}
\usepackage{parskip}
\definecolor{cdblue}{HTML}{203957}
\definecolor{cdgray}{HTML}{586677}
\hypersetup{colorlinks=true,linkcolor=cdblue,urlcolor=cdblue,pdftitle={Chapter 06: Build and Deployment Scripting}}
\pagestyle{fancy}
\fancyhf{}
\lhead{\small\color{cdgray}Continuous Delivery Engineering}
\rhead{\small\color{cdgray}Study guide}
\cfoot{\thepage\ / \pageref{LastPage}}
\setlength{\headheight}{14pt}
\setlength{\emergencystretch}{2em}
\setlist[itemize]{leftmargin=*,itemsep=3pt,topsep=4pt}
\setlist[enumerate]{leftmargin=*,itemsep=3pt,topsep=4pt}
\begin{document}
\begin{center}
{\LARGE\bfseries\color{cdblue}Chapter 06: Build and Deployment Scripting}\\[7pt]
{\large II: The Deployment Pipeline}\\[6pt]
{\small Reference edition: 2026-10-30 $\bullet$ Original educational and operational material}
\end{center}
\noindent\color{cdblue}\rule{\linewidth}{1.2pt}\color{black}
\vspace{0.4em}
\section{Chapter scope and provenance}
Primary reading: Humble and Farley, Continuous Delivery (2010), Chapter 6, "Build and Deployment Scripting"; Part II: The Deployment Pipeline. This guide is an original synthesis and not a reproduction of the chapter. The book provides the foundational delivery principles; current AWS, Terraform and GitHub Actions references in the applied sections are supplementary.


\section{Learning outcomes}
\begin{enumerate}
\item Define idempotence and deterministic execution.
\item Make a deployment script safe for repeated execution.
\item Specify meaningful exit codes and logs.
\end{enumerate}

\section{Core conceptual model}
Focus areas: idempotent scripts; repeatable deployment; build tools; environment consistency.

\begin{itemize}
\item Input: a version-controlled change, explicit requirements, and a reproducible execution context.
\item Mechanism: controlled feedback, automation, evidence capture, and a stop condition where risk exceeds tolerance.
\item Output: a verifiable release candidate or a diagnosed failure with an accountable owner.
\end{itemize}

\section{Reading and analysis protocol}
\begin{enumerate}
\item Read the chapter and record five propositions, distinguishing an author claim from an example or assumption.
\item Identify two failure modes, one recovery technique, and one measurable acceptance condition.
\item Trace each technique to the stages and controls in the deployment-pipeline-architecture reference.
\item State how the proposed practice changes for a regulated system with segregated duties.
\end{enumerate}

\section{Detailed technical study notes}
\subsection{Script boundary}
A build script should consume declared inputs and exit with a trustworthy status, not read hidden operator machine configuration.

\subsection{Idempotent deployment}
Re-running a stage should reach the same desired outcome without duplicating resources or corrupting state; record exceptions and locks.

\subsection{Build and package separation}
Compile, test, version and package in the trusted build stage; do not invent new binaries after approval.

\subsection{Cross-environment consistency}
Parameterize deployment targets instead of maintaining divergent deployment scripts for each environment.

\subsection{Portability}
Resolve paths relative to repository or module roots, and explicitly check interpreter, runtime and working-directory assumptions.

\subsection{Operational observability}
Emit clear logs, timestamps, artifact identifiers, return codes and verification results; secret values must remain masked.

\section{Domain-specific challenge questions}
\begin{enumerate}
\item For script boundary, identify the precondition, expected evidence, adversarial failure case and measurable corrective action.
\item For idempotent deployment, identify the precondition, expected evidence, adversarial failure case and measurable corrective action.
\item For build and package separation, identify the precondition, expected evidence, adversarial failure case and measurable corrective action.
\item For cross-environment consistency, identify the precondition, expected evidence, adversarial failure case and measurable corrective action.
\end{enumerate}
\section{Guided exercise}
Implement a dry-run-capable deployment script that checks prerequisites and records artifact digest.

\subsection{Required evidence}
\begin{itemize}
\item A version-controlled design or executable artifact.
\item A failure demonstration, negative test or counterexample.
\item An improvement statement with baseline, target, owner and verification method.
\end{itemize}


\section{Engineering scenario}
A deployment step succeeds when rerun manually but fails inside the pipeline due to relative paths.

\begin{itemize}
\item Describe the failure propagation path and the earliest feasible detection point.
\item Recommend a minimal remediation and its rollback or containment plan.
\item Define a measure that would demonstrate risk reduction.
\end{itemize}

\section{Knowledge checks}
\begin{enumerate}
\item Which condition is the gate, and which action is the response when it fails?
\item What evidence distinguishes successful automation from an assumed successful release?
\item Where might environments, data, dependencies or people invalidate the technique?
\item What feedback could be moved earlier without losing necessary confidence?
\item Which assumptions are book principles, and which are modern supplementary implementation choices?
\end{enumerate}

\section{Completion rubric}
Score each dimension from 0 (missing) to 4 (demonstrated and independently verifiable): conceptual accuracy; reproducibility; negative testing; security and safety; documentation and traceability. Pass threshold: 16/20 with no zero on reproducibility or safety.


\section{Prerequisites and onward route}
Chapters 1--4 and prior dependency stages. Next: Chapters 7 onward.



\end{document}
